%
% Anwendungen komplexer Funktionen und Abbildungen
%
\subsubsection{Anwendungen komplexer Funktionen und Abbildungen}

%
% elektro:
%
Real- und Imaginärteil einer komplexen Funktion sind automatisch
harmonisch.
Setzt man eine harmonische Funktion auf $\mathbb{C}$ mit einer
holomorphen Funktion zusammen, bleibt sie harmonisch.
Holomorphe Funktionen ermöglichen also, harmonische Funktionen
effizient zu manipulieren.
Die Elektrostatik beschreibt elektrische Phänomene mit einer
harmonischen Funktion.
Wie die Eigenschaften komplexer Abbildungen die Lösung
elektrostatischer Probleme ermöglichen zeigen
\emph{Malenka Hossli}
\index{Malenka Hossli}%
\index{Hossli, Malenka}%
und
\emph{Andrea Studer}
\index{Andrea Studer}%
\index{Studer, Andrea}%
in Kapitel~\ref{chapter:elektro}.

%
% joukowski:
%
Auch die zweidmensionale laminare Strömung kann mit einer harmonischen
Strömungsfunktion modelliert werden.
\emph{Jack Dumovich}
\index{Jack Dumovich}%
\index{Dumovich, Jack}%
und
\emph{Roman Meyer}
\index{Roman Meyer}%
\index{Meyer, Roman}%
zeigen in Kapitel~\ref{chapter:joukowski}, wie dies auf das Konzept
der Zirkulation und eine auf Joukowski zurückgehende Formel für den
aerodynamischen Auftrieb führt.

%
% aircraft:
%
Komplexe Zahlen können auch zur Modellierung des dynamischen
Verhaltens eines Flugzeugs herangezogen werden.
\emph{Stephanie Märklin}
\index{Stephanie Märklin}%
\index{Märklin! Stephanie}%
beschreibt in Kapitel~\ref{chapter:aircraft}, wie ein stabil fliegendes
Flugzeug in der Nähe der Ruhelage durch eine einzige Matrix charakterisiert
werden kann.
Aus den komplexen Eigenwerten dieser Matrix lässt sich ableiten, wie
stark ein Flugzeug zu unerwünschten Oszillationen wie der \emph{dutch roll}
neigt.

%
% fresnel:
%
Um die Erreichbarkeit durch Radiowellen sicherzustellen, muss man
in der drahtlosen Kommunikation verstehen, wie stark elektromagnetische
Wellen durch Hindernisse abgeschwächt werden.
Die Anwendung des huygensschen Prinzips führt auf Integrale, die sich
im Komplexen besonders leicht berechnen lassen.
\emph{Andri Sprecher}
\index{Andri Sprecher}%
\index{Sprecher, Andri}%
und
\emph{Fabian Suter}
\index{Fabian Suter}%
\index{Suter, Fabian}%
rechnen dies in Kapitel~\ref{chapter:fresnel} nach.

%
% pade:
%
Die Taylor-Approximation wird oft für die approximative Berechnung 
von Funktionen verwendet.
Die Taylor-Polynome wachsen aber unbeschränkt an, was zum Beispiel
für Funktionen wie $f(x)=e^{-x^2}$ von vornherein nicht funktionieren
kann.
Die Padé-Approximation verwendet rationale Funktionen als
Approximationsfunktionen.
Diese sind ähnlich einfach zu berechnen wie Polynome, können aber auch
beschränkte Funktionen oder Funktionen, die für $x\to\pm\infty$
konvergieren, vernünftig approximieren.
Das Kapitel~\ref{chapter:pade} von \emph{José Schmid} zeigt, wie dies geht.
\index{José Schmid}%
\index{Schmid, José}%

%
% step: Complex Step Derivative
%
In der reellen numerischen Analysis ist die Bestimmung zuverlässiger
Werte für die Ableitung ein besonderes Problem.
Ursache dafür sind die unterschiedlichen Grössenordnungen von Koordinaten
und Koordinatendifferenzen.
\index{finite Differenz}%
\index{Differenz, finite}%
Bei der Berechnung der Ableitung mit finiten Differenzen wird die
Genauigkeit durch starke Auslöschung reduziert.
\index{Ausloschung@Auslöschung}%
\emph{Selina Malacarne} zeigt in Kapitel~\ref{chapter:step}, wie die
\index{Selina Malacarne}%
\index{Malacarne, Selina}%
Ableitung in imaginärer Richtung für holomorphe Funktionen auf natürliche
Art die Trennung von Koordinaten und Koordinatendifferenz ermöglicht
und so das Problem der Auslöschung umgeht.
Eingeführt wird die Idee mithilfe der dualen Zahlen, die als eine
\index{duale Zahlen}%
Vorstufe der Methoden der Nichtstandard-Analysis zur Berechnung der
\index{Nichtstandard-Analysis}%
Ableitung angesehen werden können.

